\documentclass{ximera}
%verbeterd 12/11/26

%   \author{}
    \title{Bibliografie}
\begin{document}
\begin{abstract}
    Lijst van belangrijkste gebruikte bronnen.
\end{abstract}
\maketitle
    

\begin{itemize}
    \item Thomson Learning International Edition, \textit{Physics for scientists and engineers with Modern Physics}, Seventh edition, R. A. Serway and J.W. Jewett.
 
    \item Interactie 5.2, Die Keure
\end{itemize}
 
 
 

\section*{Bronnen van afbeeldingen}
Figuren zonder bronvermelding zijn eigen tekeningen (TikZ) van de auteur. Foto's komen van Wikimedia Commons en worden gebruikt onder de vermelde licentie: publiek domein, of een Creative Commons-licentie (CC BY of CC BY-SA, zie \url{https://creativecommons.org/licenses/}).

\subsection*{Elektrische krachtwerking}
\begin{itemize}
        \item Figuur~\verwijs[ESD-waarschuwingslabel]{fig:esd_label}{elektrische_krachtwerking_inleiding}: Tpdwkouaa, CC BY-SA 4.0, Wikimedia Commons, \url{https://commons.wikimedia.org/wiki/File:ESD_label.jpg}
        \item Figuur~\verwijs[Ballonnen trekken aan en stoten af]{fig:ballonnen}{elektrische_krachtwerking_lading}: MikeRun, CC BY-SA 4.0, Wikimedia Commons, \url{https://commons.wikimedia.org/wiki/File:Attractive-and-repulsive-electric-force-demonstrations.jpg}
        \item Figuur~\verwijs[Kat met piepschuimbolletjes]{fig:kat_piepschuim}{elektrische_krachtwerking_tribo-elekrisch_effect}: Sean McGrath, CC BY 2.0, Wikimedia Commons, \url{https://commons.wikimedia.org/wiki/File:Cat_demonstrating_static_cling_with_styrofoam_peanuts.jpg}
        \item Figuur~\verwijs[Portret van Coulomb]{fig:coulomb_portret}{elektrische_krachtwerking_coulombkracht}: Louis Hierle, publiek domein, Wikimedia Commons, \url{https://commons.wikimedia.org/wiki/File:Charles_de_Coulomb.png}
        \item Figuur~\verwijs[Torsiebalans van Coulomb]{fig:torsiebalans}{elektrische_krachtwerking_coulombkracht}: Charles-Augustin de Coulomb (1785), publiek domein, Wikimedia Commons, \url{https://commons.wikimedia.org/wiki/File:Balance_de_Coulomb.jpg}
        \item Figuur~\verwijs[Geel-groene aardingsdraad]{fig:aarding}{elektrische_krachtwerking_materialen_en_geleiding}: Erik Wannee, CC0, Wikimedia Commons, \url{https://commons.wikimedia.org/wiki/File:Installatiedraad.jpg}
        \item Figuur~\verwijs[Vandegraaffgenerator]{fig:vandegraaffgenerator_foto}{elektrische_krachtwerking_materialen_en_geleiding}: Jared C. Benedict, CC BY-SA 3.0, Wikimedia Commons, \url{https://commons.wikimedia.org/wiki/File:Van_de_graaff_generator_sm.jpg}
        \item Figuur~\verwijs[Historische bladgoudelektroscoop]{fig:elektroscoop_foto}{elektrische_krachtwerking_elektroscoop}: Museo Nazionale Scienza e Tecnologia Leonardo da Vinci, CC BY-SA 4.0, Wikimedia Commons, \url{https://commons.wikimedia.org/wiki/File:Elettroscopio_a_foglia_d'oro_-_Museo_scienza_tecnologia_Milano_11147.jpg}
        \item Figuur~\verwijs[Papiersnippers aangetrokken door een geladen cd (b)]{fig:scheiden van lading}{elektrische_krachtwerking_aantrekken_ongeladen_deeltjes}: Pediadeep, publiek domein, Wikimedia Commons, \url{https://commons.wikimedia.org/wiki/File:Paper_shavings_attracted_by_charged_cd.jpg}
        \item Figuur~\verwijs[Stofwissers met microvezels]{fig:swiffer}{elektrische_krachtwerking_aantrekken_ongeladen_deeltjes}: Lusi Lindwurm, CC BY-SA 4.0, Wikimedia Commons, \url{https://commons.wikimedia.org/wiki/File:NL21-167_microfiber-feather-dusters.jpg}
        \item Figuur~\verwijs[Poedercoating]{fig:poedercoating}{elektrische_krachtwerking_aantrekken_ongeladen_deeltjes}: Hardcoreraveman, publiek domein, Wikimedia Commons, \url{https://commons.wikimedia.org/wiki/File:Powder_Coating_Spraying.JPG}
        \item Figuur~\verwijs[Binnenkant van een elektrostatisch filter]{fig:stoffiltratie}{elektrische_krachtwerking_aantrekken_ongeladen_deeltjes}: LukaszKatlewa, CC BY-SA 3.0, Wikimedia Commons, \url{https://commons.wikimedia.org/wiki/File:Inside_of_the_electrostatic_precipitator.jpg}
        \item Figuur~\verwijs[Drum van een laserprinter]{fig:printer}{elektrische_krachtwerking_aantrekken_ongeladen_deeltjes}: Raimond Spekking, CC BY-SA 4.0, Wikimedia Commons, \url{https://commons.wikimedia.org/wiki/File:Xerox_WorkCentre_6605_-_photoconductor_drum-8998.jpg}
\end{itemize}

\subsection*{Het elektrisch veld}
\begin{itemize}
        \item Figuur~\verwijs[Michael Faraday]{fig:faraday_portret}{elektrischveld_inleiding}: Thomas Phillips (1842), publiek domein, Wikimedia Commons, \url{https://commons.wikimedia.org/wiki/File:M_Faraday_Th_Phillips_oil_1842.jpg}
        \item Figuur~\verwijs[Griesmeel toont het dipoolveld]{fig:griesmeel_dipool}{elektrischveld_dipool}: MikeRun, CC BY-SA 4.0, Wikimedia Commons, \url{https://commons.wikimedia.org/wiki/File:Semolina-grains-on-castor-oil-dipole-field.jpg}
        \item Figuur~\verwijs[Griesmeel toont schermwerking]{fig:griesmeel_schermwerking}{elektrischveld_materie_in_veld}: MikeRun, CC BY-SA 4.0, Wikimedia Commons, \url{https://commons.wikimedia.org/wiki/File:Semolina-grains-on-castor-oil-field-cancellation-inside.jpg}
        \item Figuur~\verwijs[Coaxkabel]{vb:coaxkabel}{elektrischveld_materie_in_veld}: Arj, CC BY-SA 3.0, Wikimedia Commons, \url{https://commons.wikimedia.org/wiki/File:RG-59.jpg}
        \item Figuur~\verwijs[Bliksem op een auto]{fig:bliksem_auto}{elektrischveld_materie_in_veld}: Sunzi safari, CC BY-SA 3.0, Wikimedia Commons, \url{https://commons.wikimedia.org/wiki/File:Lightning_to_a_car.jpg}
        \item Figuur~\verwijs[Gel na elektroforese]{fig:gel_foto}{elektrischveld_toepassingen}: Sgroey, CC BY-SA 4.0, Wikimedia Commons, \url{https://commons.wikimedia.org/wiki/File:GelDoc_DNA_gel_electrophoresis_photograph_stained_with_ethidium_bromide_01.jpg}
        \item Figuur~\verwijs[Vinger op een aanraakscherm]{fig:aanraakscherm_vinger}{elektrischveld_toepassingen}: Japanexperterna.se, CC BY-SA 2.0, Wikimedia Commons, \url{https://commons.wikimedia.org/wiki/File:Scrolling_on_phone_(cropped).jpg}
\end{itemize}

\subsection*{Potentiële energie, potentiaal en spanning}
\begin{itemize}
        \item Figuur~\verwijs[Bliksem]{fig:bliksem}{potentiaal_inleiding}: Hansueli Krapf, CC BY-SA 3.0, Wikimedia Commons, \url{https://commons.wikimedia.org/wiki/File:Lightning_14.07.2009_20-42-33.JPG}
        \item Figuur~\verwijs[Batterijen]{fig:batterijen}{potentiaal_spanning}: Lead holder, CC BY-SA 3.0, Wikimedia Commons, \url{https://commons.wikimedia.org/wiki/File:Batteries_comparison_4,5_D_C_AA_AAA_AAAA_A23_9V_CR2032_LR44_matchstick-1.jpeg}
        \item Figuur~\verwijs[Spanning meten met een multimeter]{fig:spanning_meten}{potentiaal_spanning}: Jiazheng0609, CC BY-SA 4.0, Wikimedia Commons, \url{https://commons.wikimedia.org/wiki/File:Testing_voltage_of_a_battery_using_a_digital_multimeter.jpg}
        \item Figuur~\verwijs[Opstelling van Millikan]{fig:millikan_opstelling}{potentiaal_energie_arbeid}: onbekende fotograaf, publiek domein, Wikimedia Commons, \url{https://commons.wikimedia.org/w/index.php?curid=685647}
        \item Figuur~\verwijs[Bliksemafleider]{fig:bliksemafleider}{potentiaal_equipotentiaaloppervlakken}: Infrogmation of New Orleans, CC BY-SA 4.0, Wikimedia Commons, \url{https://commons.wikimedia.org/wiki/File:Lightning_rod,_Woodland_Plantation_House,_LaPlace,_Louisiana.jpg}
\end{itemize}

\subsection*{Het magnetisch veld}
\begin{itemize}
        \item Figuur~\verwijs[Kompas]{fig:kompas}{magnetisme_inleiding}: Evan-Amos, publiek domein, Wikimedia Commons, \url{https://commons.wikimedia.org/wiki/File:Plastic-compass.jpg}
        \item Figuur~\verwijs[Magnetiet trekt spijkers aan]{fig:magnetiet}{magnetisme_permanente_magneten}: F. A. Annet (1922), publiek domein, Wikimedia Commons, \url{https://commons.wikimedia.org/wiki/File:Lodestone_attracting_nails.png}
        \item Figuur~\verwijs[IJzervijlsel rond een staafmagneet]{fig:ijzervijlsel_staafmagneet}{magnetisme_veld}: MikeRun en Vera Wurmsdobler, CC BY-SA 4.0, Wikimedia Commons, \url{https://commons.wikimedia.org/wiki/File:Bar-magnet-iron-filings_max.jpg}
        \item Figuur~\verwijs[IJzervijlsel tussen twee staafmagneten (a)]{fig:ijzervijlsel_twee_magneten}{magnetisme_veld}: A. W. Duff (1908), publiek domein, Wikimedia Commons, \url{https://commons.wikimedia.org/wiki/File:Magnetic_field_of_bar_magnets_attracting.png}
        \item Figuur~\verwijs[IJzervijlsel tussen twee staafmagneten (b)]{fig:ijzervijlsel_twee_magneten}{magnetisme_veld}: A. W. Duff (1908), publiek domein, Wikimedia Commons, \url{https://commons.wikimedia.org/wiki/File:Magnetic_field_of_bar_magnets_repelling.png}
        \item Figuur~\verwijs[Schooltoestel voor de proef van Oersted]{fig:proef_oersted}{magnetisme_rechte_geleider}: A. Privat-Deschanel (19de eeuw), publiek domein, Wikimedia Commons, \url{https://commons.wikimedia.org/wiki/File:Oersted_experiment.png}
        \item Figuur~\verwijs[IJzervijlsel rond een stroomvoerende draad]{fig:ijzervijlsel_rechte_geleider}{magnetisme_rechte_geleider}: Maciej J. Mrowinski, CC BY-SA 4.0, Wikimedia Commons, \url{https://commons.wikimedia.org/wiki/File:Magnetic_field_around_wire.jpg}
        \item Figuur~\verwijs[IJzervijlsel rond een stroomvoerende spoel]{fig:ijzervijlsel_spoel}{magnetisme_spoel}: Maciej J. Mrowinski, CC BY-SA 4.0, Wikimedia Commons, \url{https://commons.wikimedia.org/wiki/File:Magnetic_field_around_solenoid.jpg}
        \item Figuur~\verwijs[Hijsmagneet]{fig:hijsmagneet}{magnetisme_elektromagneet}: C. M. Jansky (1914), publiek domein, Wikimedia Commons, \url{https://commons.wikimedia.org/wiki/File:Industrial_lifting_magnet.jpg}
        \item Figuur~\verwijs[MRI-scanner]{fig:mri_scanner}{magnetisme_elektromagneet}: Jan Ainali, CC BY 3.0, Wikimedia Commons, \url{https://commons.wikimedia.org/wiki/File:MRI-Philips.JPG}
\end{itemize}

\subsection*{Appendix}
\begin{itemize}
        \item Figuur~\verwijs[Rechterhandregel]{fig:rechterhandregel}{appendix_vectoren_vectorieel_product}: Acdx, CC BY-SA 3.0, Wikimedia Commons, \url{https://commons.wikimedia.org/wiki/File:Right_hand_rule_cross_product.svg}
        \item Figuur~\verwijs[Corgi]{fig:corgi}{appendix_standaardvorm_inleiding}: Dog breed facts, CC BY-SA 4.0, Wikimedia Commons, \url{https://commons.wikimedia.org/wiki/File:Welsh_Pembroke_Corgi.jpg}
\end{itemize}

\subsection*{Titelblad}
\begin{itemize}
        \item Foto van de Joint European Torus: EUROfusion, CC BY 4.0, Wikimedia Commons, \url{https://commons.wikimedia.org/wiki/File:JET_vessel_internal_view.jpg}
\end{itemize}

\end{document}